\section*{Chapter \ref{ch:ll:websocket-rest-tls-and-json-at-speed} --- WebSocket, REST, TLS and JSON at Speed}

\begin{solution}{exo:ll:websocket-rest-tls-and-json-at-speed:1}
$2 + 2 + 4 + 300 = 308$ bytes (header, 16-bit length, masking key, payload); $2 + 8 + 4 + 70\,000 = 70\,014$ bytes; a server
frame is not masked, so $2 + 100 = 102$ bytes.
\end{solution}

\begin{solution}{exo:ll:websocket-rest-tls-and-json-at-speed:2}
Two round trips (TCP, then TLS) and the handshake's computation: $2 \times 0.5 + 2.06 \approx \qty{3.1}{\milli\second}$ with a
full handshake, about \qty{2.6}{\milli\second} with a resumed one. On the measured laptop the handshake's computation, which
both ends share, costs more than the round trips.
\end{solution}

\begin{solution}{exo:ll:websocket-rest-tls-and-json-at-speed:3}
Zero-round-trip data can be replayed: RFC~8446 buys the saved round trip ``at the cost of certain security properties''. An
order replayed by an attacker or a retransmitting middlebox would be executed twice; orders must travel only after the
handshake completes.
\end{solution}

\begin{solution}{exo:ll:websocket-rest-tls-and-json-at-speed:4}
Each padded key is exactly one 64-byte block that depends only on the key, so the hash state after it is the same for every
message. Without precomputation: the inner hash takes the ipad block and the 119-byte message padded to two blocks (three
compressions), the outer the opad block and the 32-byte inner digest padded to one (two): five. With it: three.
\end{solution}

\begin{solution}{exo:ll:websocket-rest-tls-and-json-at-speed:5}
6\,412\,320\,000\,000 exactly. In binary floating point 64123.2 is not representable; multiplied by $10^8$ it lands a hair
below or above the integer, so truncation can give 6\,412\,319\,999\,999, and prices that differ in the last digit can compare
equal or unequal by accident. Parsing the decimal text into an integer avoids the question.
\end{solution}

\begin{solution}{exo:ll:websocket-rest-tls-and-json-at-speed:6}
Answer the ping at once with a pong carrying the same payload, keep the first fragment, and deliver the depth update only
when its final fragment arrives; control frames may interleave with fragments but never split them.
\end{solution}

\begin{solution}{exo:ll:websocket-rest-tls-and-json-at-speed:7}
Repeat the key to the payload's length and turn both byte strings into integers
(\texttt{int.from\_bytes}); XOR the two integers and convert the result back to bytes (\texttt{to\_bytes}). The fixture tests are unchanged; the cost per 300-byte frame falls
from tens of microseconds to a few.
\end{solution}

\begin{solution}{exo:ll:websocket-rest-tls-and-json-at-speed:8}
Every order pays two round trips and a handshake, about \qty{3.1}{\milli\second} here, the venue sees a new connection per
order (and may limit or ban them), and the rate-limit governor has no single place to count requests. A pool of \omterm{def:ll:websocket-rest-tls-and-json-at-speed:rest}{persistent
connections} shared behind the governor avoids all three; contention for a connection costs nanoseconds, a handshake
milliseconds.
\end{solution}

\begin{solution}{pb:ll:websocket-rest-tls-and-json-at-speed:1}
\begin{enumerate}
\item On the committed measurement, about \qty{0.9}{\micro\second} with the C++ components and \qty{48}{\micro\second} with the
Python reference.
\item C++: the JSON decoding (\qty{0.39}{\micro\second}). Python: the byte-by-byte framing, \qty{16}{\micro\second} to unmask
and \qty{15}{\micro\second} to mask.
\item $0.93/(0.93 + 500) \approx 0.2\%$ and $47.6/(47.6 + 500) \approx 8.7\%$.
\item About \qty{16.9}{\micro\second}, $3.3\%$ of the path: framing was two-thirds of the Python cost.
\item $2 \times \qty{0.5}{\milli\second}$ of round trips plus the handshake: about \qty{3.1}{\milli\second} (full) or
\qty{2.6}{\milli\second} (resumed) per order.
\item Three round trips and the handshake, about \qty{3.6}{\milli\second}.
\item Opened one after the other, about 330 a second; on a \omterm{def:ll:websocket-rest-tls-and-json-at-speed:rest}{persistent connection} the limit is the venue's rate limit.
\item Plan a reconnection before the limit: open the new connection, subscribe and resynchronise the book, then close the old
one, so that no update is missed.
\item It covers the order's price and quantity, known only after the decision. The keyed hash state, the fixed parts of the
payload and the formatting of the timestamp can be prepared in advance.
\item 5\,000 updates a second: about 0.2\% of a core with the indexed decoder and 4.9\% with the Python reference.
\item Missed updates: a first identifier above the last applied one plus one. The book is discarded and rebuilt from a new
snapshot (Book~3's \texttt{firm.wsbook}).
\item So that a captured request cannot be replayed later: the venue refuses one older than its window.
\item \textbf{Named result.} With a \qty{250}{\micro\second} hop each way, decryption, framing, JSON and signing take 0.2\% of
the venue--host--venue path with the C++ components and 8.7\% with the Python reference; reusing a connection saves about
\qty{3.1}{\milli\second} per REST order against opening one (\qty{2.6}{\milli\second} with resumption), several times the
network path itself.
\item When the strategy's decisions are worth milliseconds, not microseconds, and messages are few: the network hop and the
venue's own latency dominate anyway.
\item In a separate thread or process that terminates TLS for several connections, or in the operating system's or the
network card's TLS support where available, so that the strategy's thread sees plain frames.
\item The same costs on the server's processor (its AES and SHA instructions), under the real message rates and sizes, and
the distribution of the network hop, not just its median.
\item It is not processor time on this host, but it is most of the tick-to-trade and its variance; a closer placement changes
the result more than any decoder.
\item JSON and decimal text, and the masking; the parsing becomes fixed-offset reads (chapter~16).
\item Feed the pool a scripted sequence of venue answers with 429 and a retry time, as Book~3's governor accepts through
\texttt{on\_status}, and check that it stops sending until the time has passed.
\item It keeps connections open, decodes without building anything and signs without rehashing the key.
\end{enumerate}
\end{solution}

\begin{solution}{iq:ll:websocket-rest-tls-and-json-at-speed:1}
A new connection costs a TCP round trip, a TLS round trip and the handshake's cryptography (milliseconds), and may count
against the venue's connection limits; a \omterm{def:ll:websocket-rest-tls-and-json-at-speed:rest}{persistent connection} pays this once. Use a pool behind the rate-limit governor.
\iqlookfor{round trips plus handshake computation, and limits.}
\end{solution}

\begin{solution}{iq:ll:websocket-rest-tls-and-json-at-speed:2}
Two header bytes (final bit, opcode, mask bit, 7-bit length), an extended 16- or 64-bit length, a 4-byte masking key on client
frames, the payload. Masking protects intercepting proxies from crafted payloads; it costs an XOR per byte, nanoseconds when
done eight or thirty-two bytes at a time.
\iqlookfor{layout, the reason, and the cost done well.}
\end{solution}

\begin{solution}{iq:ll:websocket-rest-tls-and-json-at-speed:3}
Find all structural characters in one vectorised pass, then walk that index with the message's expected shape, reading
numbers in place and parsing prices into integers; fall back to a general parser if the shape differs. No tree, no allocation,
no floating point.
\iqlookfor{a structural index and a shape-specific walk.}
\end{solution}

\begin{solution}{iq:ll:websocket-rest-tls-and-json-at-speed:4}
HMAC-SHA-256 of the query string and body with the secret key, sent in hex with the API key, a timestamp and a validity window
in the payload. Hash the key's two padded blocks once and copy those states per request; use a library SHA-256 that uses the
processor's instructions; prepare the fixed parts of the payload.
\iqlookfor{the RFC 2104 precomputation, and replay protection.}
\end{solution}

\begin{solution}{iq:ll:websocket-rest-tls-and-json-at-speed:5}
Resumption skips the certificate and part of the key exchange by using a key from an earlier connection: less computation,
still one round trip. Zero-round-trip data is sent with the first message and can be replayed, so it must be idempotent,
which an order is not.
\iqlookfor{computation saved versus round trips, and replay.}
\end{solution}

\begin{solution}{iq:ll:websocket-rest-tls-and-json-at-speed:6}
Connection setup on the order path (a new connection or a reconnection per order), queueing behind a rate limiter or a busy
connection, garbage collection or slow decoding in the client, TLS or signing done inefficiently, and the venue's own latency;
measure each step with timestamps at every boundary.
\iqlookfor{measure the steps; handshakes first.}
\end{solution}
